% Generated draft: requires mathematical and source review.
\documentclass[11pt]{article}
\usepackage[margin=25mm]{geometry}
\usepackage{fontspec}
\setmainfont{Latin Modern Roman}
\usepackage{amsmath,amssymb,mathtools,unicode-math}
\setmathfont{Latin Modern Math}
\usepackage{enumitem,needspace,xurl,hyperref}
\setlength{\parindent}{0pt}
\setlength{\parskip}{5pt}
\setlength{\emergencystretch}{3em}
\newcommand{\N}{\mathbb{N}}
\newcommand{\Z}{\mathbb{Z}}
\newcommand{\Q}{\mathbb{Q}}
\newcommand{\R}{\mathbb{R}}
\newcommand{\C}{\mathbb{C}}
\newcommand{\F}{\mathbb{F}}
\newcommand{\A}{\mathbb{A}}
\newcommand{\PP}{\mathbb{P}}
\newcommand{\E}{\mathbb{E}}
\newcommand{\eps}{\varepsilon}
\newcommand{\dd}{\,\mathrm{d}}
\begin{document}
\section*{UnsolvedMath ID: 30005655}
\textbf{Source problem code:} catalogue entry
\subsection{Definitions and mathematical statement}
Let \(G\) be a finite group. A \(G\)-Mackey functor \(M\) assigns abelian groups to finite \(G\)-sets, with restriction and transfer maps satisfying additivity and the Mackey compatibility axioms. Write \(HM\) for the genuine \(G\)-equivariant Eilenberg--MacLane spectrum whose degree-zero homotopy Mackey functor is \(M\) and whose other integer-graded homotopy Mackey functors vanish.

\textbf{Unresolved multiplicative hypothesis.} A Mackey functor alone supplies additive data. It does not specify a commutative ring spectrum structure on \(HM\). Such additional multiplicative data are needed for the category of \(HM\)-modules to have the relative symmetric tensor product used in the stated Balmer spectrum. The supplied record does not specify these data or the intended class of multiplicative Mackey functors. Consequently, its formulation for an arbitrary Mackey functor does not yet determine a tensor-triangulated category.

Conditional on specifying a commutative algebra structure on \(HM\) in genuine \(G\)-spectra, set
\[
\mathcal K=\operatorname{Perf}_G(HM)
 :=\operatorname{Ho}(\operatorname{Mod}_{HM}(\operatorname{Sp}^{G}))^{c},
\]
where the superscript \(c\) denotes compact objects: an object \(X\) is compact when \(\operatorname{Hom}(X,-)\) preserves arbitrary coproducts. Equip \(\mathcal K\) with the derived relative tensor product \(\otimes_{HM}\) and unit \(HM\).

A thick tensor ideal \(\mathcal P\subseteq\mathcal K\) is a full triangulated subcategory closed under direct summands and under tensoring with arbitrary objects of \(\mathcal K\). It is prime if it is proper and
\[
X\otimes_{HM}Y\in\mathcal P
\quad\Longrightarrow\quad
X\in\mathcal P\ \text{or}\ Y\in\mathcal P.
\]
The Balmer spectrum \(\operatorname{Spc}(\mathcal K)\) is the set of prime thick tensor ideals, with closed subsets
\[
Z(\mathcal E)=\{\mathcal P\in\operatorname{Spc}(\mathcal K):
\mathcal P\cap\mathcal E=\varnothing\},
\]
for collections \(\mathcal E\) of objects of \(\mathcal K\).

\textbf{Requested description, conditional on clarification.} For finite groups \(G\) and the intended admissible coefficients \(M\), describe the points and topology of
\[
\operatorname{Spc}\bigl(\operatorname{Perf}_G(HM)\bigr).
\]
The supplied statement quantifies over Mackey functors, whereas the supplied literature assessment concerns constant Mackey coefficients associated with arbitrary commutative rings. These scopes cannot be identified without further source clarification; no restriction to constant coefficients is imposed here.

\textbf{Supplied contextual comparisons.} The extracted source presents the nonequivariant comparison, for a commutative ring \(R\),
\[
\operatorname{Spc}(D^{\mathrm{perf}}(R))
\cong\operatorname{Spc}(\operatorname{Perf}(HR))
\cong\operatorname{Spc}((HR\text{-}\operatorname{Mod})^{c})
\cong\operatorname{Spec}(R).
\]
Here \(D^{\mathrm{perf}}(R)\) is the category of perfect complexes of \(R\)-modules, and \(\operatorname{Spec}(R)\) has the Zariski topology. The extract attributes the comparison with \(\operatorname{Spec}(R)\) to Hopkins--Neeman, references [6, 7], when \(R\) is Noetherian, and to Thomason, reference [10], for general \(R\). It also states that, for a prime \(p\), every nonequivariant \(H\F_p\)-module \(X\) admits an equivalence
\[
X\simeq\bigvee_{i\in I}\Sigma^{n_i}H\F_p,
\qquad n_i\in\Z,
\]
for an indexing set \(I\). These are supplied background statements, not solutions of the equivariant description problem.
\subsection{Short English statement}
Describe the prime thick tensor ideals and their topology for compact modules over an equivariant Eilenberg--MacLane spectrum. The statement requires clarification: an arbitrary Mackey functor does not specify the commutative multiplication needed for this tensor product, and the supplied literature assessment addresses constant coefficients from commutative rings. This is a statement-only draft requiring human mathematical review.
\subsection{Sources}
\begin{itemize}
\item Local UnsolvedMath registry, supplied complete record: permanent ID \(30005655\), source problem code \texttt{catalogue entry}, title \emph{Balmer Spectra of Equivariant Eilenberg--MacLane Modules}. The record lists proposed year 2023 and supplies no named proposer.
\item Supplied source citation: \emph{Tensor-Triangular Geometry and Interactions} (2024), Oberwolfach Reports, pp. 2303--2358, DOI \texttt{10.4171/owr/2023/40}. \url{https://doi.org/10.4171/owr/2023/40}. The supplied background describes these reports as records of mathematical workshops and their open-problem discussions.
\item Supplied literature source: \emph{Quillen stratification in equivariant homotopy theory}. \url{https://doi.org/10.1007/s00222-024-01301-0}.
\item Supplied literature source: OWR report. \url{https://publications.mfo.de/bitstream/handle/mfo/4102/OWR_2023_40.pdf?sequence=4}.
\item Supplied contextual attributions: Hopkins--Neeman [6, 7] for the Noetherian comparison and Thomason [10] for general commutative rings. The supplied extract does not include the corresponding bibliography entries.
\item Literature assessment dated exactly \texttt{2026-08-22}: the record labels the problem open and states that the constant Mackey Eilenberg--MacLane case has computations for important field and integral coefficients, while a general description for arbitrary commutative rings remains open. This is supplied dated curation, not an independently verified literature assessment; its status and cited claims require independent verification.
\item The registry labels the statement \texttt{corrected\_verified} and reports that the official Oberwolfach report was checked directly. That verification claim is part of the supplied provenance; no source visit or independent verification is claimed in this draft.
\end{itemize}
\end{document}
